Signal retiming and congestion pricing each significantly reduce delay on a
calibrated simulation of Kathmandu's arterial network, and combining them produces
the largest effect of any intervention tested here (65.30\% $\pm$ 4.47\%). We
tested directly whether this combined effect is super-additive and found it is
not -- a correction to this project's own earlier working conclusion, reported
here rather than carried forward uncorrected. A separate robustness check across a
range of per-lane road capacities, since this study's own capacity assumption was
never measured against the city, finds the combined lever lowers delay in
direction at every capacity tested while completing more trips, whereas the
individual levers' effects are not robust below the calibrated default, meaning the standalone findings should be read
as conditional on that assumption and the combined result as this paper's more
robust claim. We report this result alongside the
honest limits of the model it rests on: a calibration that matches the city on
average but not every intersection, a route-choice process rather than a settled
equilibrium, and two further candidate interventions that we tested and rejected
rather than omitted. Two additional interventions -- staggered departure
scheduling and decentralized trip attraction -- show comparably large effects and
are reported with their own stated confounds rather than as additional headline
results. Future work should prioritize closing the validation gap with more
observed intersection counts, isolating the confounds on the two supported
additional interventions, running more seeds on the interaction term itself,
since the current result (CI including zero, 9/15 seeds negative) rules out a
large super-additive effect but cannot rule out a smaller one at this sample size,
and re-testing the individual levers below the calibrated capacity with a metric
that counts unfinished trips.
